%!TEX root = ../../main.tex
\section{본 논문의 접근과 주요 결과}
\label{sec:intro-approach}

본 논문은 교전 사례의 구성, 결과의 정의, 사전 예측, 적용 범위의 확인이라는 네 단계로 이루어진다.

교전은 공개 기록에 남은 킬로 나눈다. 서로 가까운 킬을 한 교전으로 묶을 때 허용하는 시간 간격과, 그 킬들이 지도에서 떨어져 있어도 되는 거리는 모은 경기에서 추정한다. 양 팀이 첫 킬 직전에 그 자리 근처에 있었는지를 보는 거리와 시간은 게임이 쓰는 규칙의 값을 가져온다. 이렇게 찾은 교전은 앞 절의 기준대로 한타와 소규모 교전으로 나눈다. 예측의 대상은 교전이고, 비교는 이 두 집단 안에서 따로 한다.

첫째 모형은 경기의 승패로 학습한 뒤 고정하고, 앞 절에서 정한 예측 시점과 결과 시점의 상태에 각각 승률을 매긴다. 그 증가 여부가 교전의 결과이다. 첫째 모형을 학습하는 데 쓰인 경기의 교전은 그 경기를 빼고 다시 만든 모형이 매기고, 그 밖의 경기는 고정한 모형이 매긴다. 둘째 모형은 이 증가 여부를 예측하도록 첫째 모형 위에 설계한다. 입력은 예측 시점까지 공개된 상태이고, 한타와 소규모 교전은 같은 구조로 따로 학습한다. 비교 상대는 시작 승률과 경기 시각만 쓰는 기준선이며, 같은 교전끼리 견준다. 한 경기의 교전은 한데 묶어 다시 뽑아 그 차이가 놓이는 범위를 구한다. 학습과 설정 선정에 쓰는 경기와 마지막으로 점수를 내는 경기는 나눈다. 다만 점수를 내는 시기의 경기를 절차를 굳히기 전에 살펴본 이력이 있으므로, 주 비교는 한 번도 보지 않은 자료에서 확정한 결과가 아니라 탐색적으로 읽는다.

틀린 정도는 확률 예측의 오차인 Brier 점수로 재며, 낮을수록 덜 틀린 것이다. 한타와 소규모 교전 모두에서 둘째 모형은 기준선보다 덜 틀렸고, 경기를 다시 뽑아도 그 방향은 유지된다. 둘째 모형이 기준선에 더하는 개선은 두 코호트 모두 절대적으로 작다. 다만 그 작음이 같은 뜻은 아니다. 한타에서는 모든 교전에 오를 확률이 절반이라고만 답하는 것보다 나아진 부분의 상당수를 시작 승률과 경기 시각이 차지한다. 소규모 교전에서는 그 기준선이 절반이라고만 답하는 것보다 나은 폭이 제한적이고, 나머지 사전 상태를 넣었을 때 추가 개선이 나타난다.

어느 쪽도 크게 앞서 있지 않은 교전에서도 사전 상태가 더하는 작은 개선은 관측된다. 한타에서는 그 개선을 다시 뽑아 구한 범위가 차이가 없는 쪽에 가까이 가서, 뚜렷한 크기의 이득이라고 말하기는 어렵다. 나중 시기와 다른 지역에는 모형을 다시 학습하지 않고 점수만 매겼다. 한타에서는 둘째 모형이 기준선보다 나은 경우가 없었다. 소규모 교전에서는 조금 나았으나, 지역과 코호트를 한 묶음으로 보고 기준을 높이면 그 우세는 확실하지 않다.

이 결과는 고정해 둔 첫째 모형, 교전을 나눈 방식, 실제로 평가한 절차와 자료 안에서 읽는다. 사람이 표시한 교전의 시작과 끝, 그리고 그 교전이 승리를 일으켰다는 인과는 이 논문이 다루는 대상이 아니다. 수치는 제 \ref{ch:engagement}--\ref{ch:validation} 장에서 제시한다.
